\section{Method III: From personalization to post-training (\fsppo)}
\label{sec:post-training}

The paper does not stop at $\widehat Q_H$. We update the weights of an open-weight language model.
Let $\pi_{\mathrm{ref}}$ be the initial policy and $\pi_\vartheta$ the policy being trained. (We write
$\vartheta$ for policy parameters to avoid a clash with the smoothing parameter $\theta$.) For a
comparison $X$, the DPO implicit utility difference \citep{rafailov2023dpo} is
\[
  d_\vartheta(X)=\beta\Bigl[\log\frac{\pi_\vartheta(y_A\mid x)}{\pi_{\mathrm{ref}}(y_A\mid x)}
  -\log\frac{\pi_\vartheta(y_B\mid x)}{\pi_{\mathrm{ref}}(y_B\mid x)}\Bigr].
\]

\paragraph{Binary human labels.}
With $\hat p_{H,i}=\hmu(Z_i)$ from \dfsp{}, the soft-label objective is
\[
  \widehat R(\vartheta;\hat p_H)=\frac1N\sum_{i=1}^{N}\ell\bigl(d_\vartheta(X_i);\hat p_{H,i}\bigr),\qquad
  \ell(d;p)=-p\log\sigma(d)-(1-p)\log\sigma(-d).
\]
This loss is deliberately standard; soft-label DPO and GAPO already use soft targets
\citep{lee2024rlaif,furuta2024gapo}. The contribution is where $\hat p_H$ comes from and what
guarantee it carries. \dfsp{} targets can also be plugged into GAPO's objective. We test that
empirically, but the guarantees below cover only the soft cross-entropy loss.

\paragraph{Bounded scorer.}
All guarantees assume $|d_\vartheta|\le D$. We treat this as a \emph{constraint on the analyzed
class}. In training it is implemented by clipping $d_\vartheta$ inside the loss, so the analyzed
scorer is the clipped one that actually enters the objective. Proposition~\ref{prop:qg} needs a
convex class, and therefore the constraint form rather than clipping.

\paragraph{Ordinal or richer labels.}
With a cumulative-link model \citep{mccullagh1980regression},
$P_\vartheta(S=s\mid X)=G(\tau_{s+1}-d_\vartheta)-G(\tau_s-d_\vartheta)$, we minimize
$-\sum_i\sum_s\widehat Q_H(s\mid X_i)\log P_\vartheta(S=s\mid X_i)$. Ordinal and strength-aware
preference optimization already exist \citep{imai2026mixdpo}; what is new is the personalized target.

\subsection{Bridge: personalization error controls post-training excess risk}

Let $R(\vartheta;p)=\E_X\,\ell(d_\vartheta(X);p(X))$, let $p_X=P(Y=1\mid X)$, and recall
$R_X(\hat p)=\E(\hat p-p_X)^2=R(\hat p)+B_\zeta$. All three results below rest on one identity: since
$\log\sigma(d)-\log\sigma(-d)=d$, the loss is \emph{linear in the target},
$\ell(d;p)-\ell(d;p')=-(p-p')\,d$. They are standard target-perturbation and strong-convexity
arguments. We state them because they turn the personalization rate into a post-training rate.

\begin{proposition}[General bridge]
\label{prop:bridge}
Assume $|d_\vartheta|\le D$. Let $\hat p$ be fixed (for example, cross-fitted on labels independent
of the policy-optimization sample). Let
$R(\hat\vartheta;\hat p)\le\inf_\vartheta R(\vartheta;\hat p)+\varepsilon$, where $\varepsilon$ collects
generalization error over the $N$ pairs and optimization error. Then
\[
  R(\hat\vartheta;p_X)-\inf_\vartheta R(\vartheta;p_X)\;\le\;\varepsilon+2D\,\E|\hat p(X)-p_X|
  \;\le\;\varepsilon+2D\,R_X(\hat p)^{1/2}.
\]
\end{proposition}

\begin{proof}
By linearity, $|R(\vartheta;\hat p)-R(\vartheta;p_X)|\le D\,\E|\hat p-p_X|$ for every $\vartheta$. For
any $\vartheta'$, write $R(\hat\vartheta;p_X)-R(\vartheta';p_X)$ as
$[R(\hat\vartheta;p_X)-R(\hat\vartheta;\hat p)]+[R(\hat\vartheta;\hat p)-R(\vartheta';\hat p)]
+[R(\vartheta';\hat p)-R(\vartheta';p_X)]$. Bound the first and third brackets by linearity and the
second by $\varepsilon$, then take the infimum over $\vartheta'$. The last step is Cauchy--Schwarz.
\end{proof}

\begin{proposition}[Fast rate under quadratic growth; no realizability needed]
\label{prop:qg}
Suppose the induced scorer class $\mathcal D=\{d_\vartheta\}$ is a convex, closed subset of
$L^2(P_X)$, with $\norm{d}_\infty\le D$ for all $d\in\mathcal D$. Let $d^{*}$ minimize $R(\cdot;p_X)$
over $\mathcal D$, and let $c_D=\sigma(D)\{1-\sigma(D)\}$. If
$R(\hat d;\hat p)\le\inf_{\mathcal D}R(\cdot;\hat p)+\varepsilon$, then
\[
  R(\hat d;p_X)-R(d^{*};p_X)\;\le\;2\varepsilon+\frac{2}{c_D}\,R_X(\hat p).
\]
\end{proposition}

\begin{proof}
Let $\mathrm{Exc}=R(\hat d;p_X)-R(d^{*};p_X)$ and $a=\norm{\hat p-p_X}_{L^2}$. Linearity and the
near-optimality of $\hat d$ under $\hat p$ give
$\mathrm{Exc}\le\varepsilon+\E\{(\hat p-p_X)(\hat d-d^{*})\}\le\varepsilon+a\norm{\hat d-d^{*}}$.

For quadratic growth, note that $\partial_d^2\ell(d;p)=\sigma(d)\{1-\sigma(d)\}\ge c_D$ on $[-D,D]$,
whatever $p$ is. Along the segment $d^{*}+t(\hat d-d^{*})$, which stays in $\mathcal D$, the risk
therefore has second derivative at least $c_D\norm{\hat d-d^{*}}^2$. Its first derivative at $t=0$
is $\ge0$ by optimality over the convex set. Hence $\mathrm{Exc}\ge\tfrac{c_D}2\norm{\hat d-d^{*}}^2$.

Combining the two bounds, $\mathrm{Exc}\le\varepsilon+a\sqrt{2\mathrm{Exc}/c_D}\le\varepsilon+\mathrm{Exc}/2+a^2/c_D$.
\end{proof}

Convexity of $\mathcal D$ holds for scorers that are linear in their parameters, such as fixed
features or a linearized (NTK-style) regime. For LoRA-trained policies it is an idealization. The
constant $c_D\approx e^{-D}$ degrades quickly as $D$ grows.

\begin{proposition}[Realizable case]
\label{prop:bridge-fast}
Suppose $p_X,\hat p\in[\kappa,1-\kappa]$. For $\hat p$ this can be enforced by projection, which
does not increase the error. Suppose also that every measurable $p:\operatorname{supp}(P_X)\to[\kappa,1-\kappa]$ in
the relevant class is realized by some $\vartheta$, i.e.\ $\sigma(d_\vartheta)=p$ $P_X$-almost surely,
with $D\ge\operatorname{logit}(1-\kappa)$. This is a condition on the \emph{population}, not on the
finite pool.
If $\hat\vartheta$ minimizes $R(\cdot;\hat p)$ exactly, then
\[
  R(\hat\vartheta;p_X)-\inf_\vartheta R(\vartheta;p_X)=\E\,\KL\bigl(\Ber(p_X)\,\|\,\Ber(\hat p)\bigr)
  \;\le\;\frac{R_X(\hat p)}{\kappa(1-\kappa)} .
\]
\end{proposition}

\begin{proof}
Pointwise, $\ell(d;p)-H(p)=\KL(\Ber(p)\,\|\,\Ber(\sigma(d)))$, where $H$ is the binary entropy. So
$\sigma(d_{\hat\vartheta})=\hat p$ almost surely and $\inf_\vartheta R(\vartheta;p_X)=\E H(p_X)$.
Finally, $\KL\le\chi^2=(p_X-\hat p)^2/\{\hat p(1-\hat p)\}$.
\end{proof}

\emph{When is realizability coherent?} Utility differences at a fixed prompt satisfy
$d(A,B)+d(B,C)+d(C,A)=0$. So if a prompt's comparisons contain a cycle, not every assignment of
probabilities can be realized; three cyclic probabilities of $0.75$ cannot. If each prompt's
comparison graph is a \emph{forest} (for example, one pair per prompt, as in most pairwise datasets),
any assignment of pairwise probabilities is consistent with some utility function. The assumption
then reduces to the richness of the policy class. The forest condition is an assumption on the
population comparison design, meaning the support of $P_X$ given each prompt. Checking it on the pool
(E0) is a sanity check, not a proof.

\keybox{\textbf{Chain.} $n$ human labels
$\xrightarrow{\ \text{Targets \ref{tr:upper}, \ref{tr:oracle}}\ }$
$R_X(\hat p)\lesssim r_n+B_\zeta$
$\xrightarrow{\ \text{Props.\ \ref{prop:bridge}--\ref{prop:bridge-fast}}\ }$
excess preference risk $\lesssim\varepsilon+\sqrt{r_n+B_\zeta}$ (bounded scorer), or
$\lesssim\varepsilon+r_n+B_\zeta$ (quadratic growth or realizable).
The coarsening gap $B_\zeta$ does not vanish with more human labels.}

\subsection{End-to-end corollary}

\begin{target}[End-to-end minimax rate; a corollary of Target~\ref{tr:lower} and the bridge]
\label{tr:e2e}
Consider binary labels under the following assumptions:
\begin{itemize}
  \item $Z$-sufficiency: $p_X=p_H^{*}(Z)$, so $B_\zeta=0$;
  \item the population comparison design is a forest at each prompt;
  \item the policy class realizes every utility assignment on $\operatorname{supp}(P_X)$ with
        $|d|\le D$, where $D\ge\operatorname{logit}(1-\kappa)$;
  \item the judge satisfies the slack conditions of Target~\ref{tr:lower};
  \item the policy step is idealized ($N\to\infty$, $\varepsilon=0$), which isolates the cost of
        human labels.
\end{itemize}
Admissible procedures choose the $n$ revealed labels by any pool-based design, sequential or not,
and then output a policy. Then:
\begin{enumerate}[label=(\roman*)]
  \item \emph{fixed-class minimax rate:}
        \[
          \inf\,\sup_{\mathcal F^{\kappa}(\theta^{*},\gamma)}\E\bigl[R(\hat\vartheta;p_X)-\inf_\vartheta R(\vartheta;p_X)\bigr]
          \asymp_\kappa\gamma_1^{\frac{2d}{2\gamma_2+d}}n^{-\frac{2\gamma_2}{2\gamma_2+d}}+\frac1n .
        \]
        The upper bound is
        attained by \dfsp{} with $(\theta,h)$ tuned to the class, followed by
        Proposition~\ref{prop:bridge-fast}.
  \item \emph{adaptive procedure:} with $(\hat\theta,\hat h)$ chosen on validation data, the same holds
        up to the $(\log n)^2/n$ remainder of Target~\ref{tr:oracle}. This remainder dominates when
        $\gamma_1$ is very small, so the adaptive procedure does not attain the sharp $1/n$ regime.
\end{enumerate}
\end{target}

\emph{Lower bound:} Pinsker's inequality gives excess
$=\E\,\KL(\Ber(p_X)\,\|\,\Ber(\sigma(d_{\hat\vartheta})))\ge2\,\E\{\sigma(d_{\hat\vartheta})-p_X\}^2$.
The policy therefore induces an estimator $\sigma(d_{\hat\vartheta})$ of $p_H^{*}$ from $n$ adaptively
chosen labels. Under the forest and richness conditions every member of
$\mathcal F^{\kappa}(\theta^{*},\gamma)$ is realizable, so Target~\ref{tr:lower}, which covers sequential
designs, applies without shrinking the class.

\emph{Status:} under sufficiency, realizability and an idealized policy step, the policy problem
reduces to bounded Bernoulli regression through the KL--MSE equivalence. The corollary is a useful
translation of Target~\ref{tr:lower} into post-training language. It is not an independent
resolution of realistic policy-learning difficulty, where $\varepsilon$, non-convexity and the
coarsening gap dominate.

\keybox{\textbf{\textcolor{BrickRed}{Design decision D2: data reuse.}} (a) Cross-fit: split the pool
into two folds, fit \dfsp{} on one fold's revealed labels, and produce targets on the other.
(b) On the $n$ revealed comparisons, use $\hat p$ (lower variance), the observed $Y$ (unbiased), or
a mixture. The default is (a) with $\hat p$ everywhere; (b) is an ablation.}
